\documentclass[a4paper, 11pt]{article}
\usepackage[utf8]{inputenc}
\usepackage[
  left = 20mm, right = 20mm, top = 22mm, bottom = 22mm,
]{geometry}

\usepackage{amsmath,amsfonts,amssymb,bm}
\usepackage[english]{babel}
\usepackage{natbib}
\usepackage{hyperref}
\usepackage{orcidlink}
\usepackage{xcolor}
\usepackage{fancyhdr}
\usepackage{graphicx}
\usepackage{booktabs}
\usepackage{float}

\definecolor{harvardcrimson}{HTML}{8C1D40}

\hypersetup{
  colorlinks=true,
  linkcolor=harvardcrimson,
  citecolor=harvardcrimson,
  urlcolor=harvardcrimson,
  hypertexnames=false,
}

\pagestyle{fancy}
\fancyhf{}
\rhead{\textcolor{harvardcrimson}{\textbf{\texttt{y3\_cluster\_cpp}: factoring $R$ out of the \texttt{DSigma1h} mass integral}}}
\lhead{\textcolor{harvardcrimson}{\textbf{J.\,H.\,Esteves}}}
\cfoot{\thepage}
\renewcommand{\headrulewidth}{0.4pt}
\renewcommand{\footrulewidth}{0.4pt}
\setlength{\headheight}{14pt}

\renewcommand{\baselinestretch}{1.15}
\setlength{\emergencystretch}{3em}
\sloppy

\renewcommand{\topfraction}{0.95}
\renewcommand{\bottomfraction}{0.9}
\renewcommand{\textfraction}{0.05}
\renewcommand{\floatpagefraction}{0.95}
\setlength{\intextsep}{8pt}
\setlength{\textfloatsep}{10pt}
\setlength{\floatsep}{10pt}

\newcommand{\hMpc}{h^{-1}\,\mathrm{Mpc}}
\newcommand{\hMsun}{h^{-1}\,M_{\odot}}
\newcommand{\avg}[1]{\langle #1 \rangle}
\newcommand{\lnM}{\ln M}
\newcommand{\rs}{r_s}
\newcommand{\reff}{\hat r_s}
\newcommand{\reffi}[1]{\hat r_{s,#1}}
\newcommand{\yeff}{\hat y}
\newcommand{\Aeff}{\hat A}
\newcommand{\dVdz}{\dfrac{dV}{d\Omega\,dz}}
\newcommand{\DSnfw}{\Delta\Sigma_{\rm NFW}}
\newcommand{\rmis}{R_{\rm mis}}
\newcommand{\DSoh}{\Delta\Sigma_{1h}}
\newcommand{\obsij}[1]{\langle \DSoh(#1)\rangle_{ij}}
\newcommand{\Wcal}{\mathcal{W}}
\newcommand{\Lop}{\mathsf{L}}
\newcommand{\pyfile}[1]{{\small\texttt{#1}}}

\title{\texorpdfstring{%
\textcolor{harvardcrimson}{\textbf{\texttt{y3\_cluster\_cpp}:\\
factoring the radial dependence out of the \texttt{DSigma1h} mass integral}}%
}{y3_cluster_cpp: factoring the radial dependence out of the DSigma1h mass integral}}
\author{Johnny H.\ Esteves\,\orcidlink{0000-0003-4373-2386}\\[2pt]
{\small Department of Physics, Harvard University}}
\date{\today}

\begin{document}
\maketitle
\thispagestyle{fancy}

\paragraph{What this note is about, in one paragraph.} The pipeline
predicts the stacked 1-halo excess surface density,
$\obsij{R}$, for every richness bin $i$ and redshift bin $j$, by
running a fresh 2-D quadrature over halo mass and redshift for
\emph{every single radius $R$ on the grid}. That is the slow part
of the whole lensing calculation: cost grows as
$N_{\rm radii}\times N_{\rm bins}\times(\text{mass--redshift
quadrature cost})$, and the radius grid is the multiplier we would
most like to get rid of. This note walks through two separate
ideas for doing that, one exact and one approximate, building each
step from a plain-language argument before writing down the
formula. \S\ref{sec:exact} shows that the redshift integral can be
removed \emph{completely for free} --- no approximation, just a
reordering of the integral --- because of a specific, checkable
fact about this model's NFW profile. \S\ref{sec:setup}--\S\ref{sec:approaches}
then tackle the harder part, the radial dependence hiding inside
the remaining mass integral, and lay out three increasingly
sophisticated ways to pull it out. \S\ref{sec:comparison} compares
them side by side and recommends where to start.
\S\ref{sec:validation} then checks all of this numerically: the NFW
shape function is cross-checked against an independent projection of
the 3-D density, and the three ideas are tested against the exact
mass integral for a representative mass weight spanning the
$\sigma_{\lnM}$ range discussed throughout. \S\ref{sec:mis} then
extends the same machinery to the pipeline's actual observable ---
the miscentered mixture, \texttt{DSigma1hMisSel} --- and shows this
needs no new derivation, only a small generalisation of Idea 2.
Finally, \S\ref{sec:real_validation} repeats the whole check against
a real running \texttt{cosmosis} sample --- real mass function, real
selection function, real halo-model and miscentering tables, cross-checked
directly against \texttt{NumCountsSel.so} and \texttt{Shear1hMisSel.so}
--- confirming the synthetic-weight conclusion holds with the real
pipeline ingredients.

\tableofcontents

\section{The observable, and why it is expensive}
\label{sec:exact}

\subsection{What \texorpdfstring{$\obsij{R}$}{<DeltaSigma\_1h(R)>\_ij} means, in words}

Picture a bin of galaxy clusters, selected because they have
observed richness in some range $i$ and photometric redshift in
some range $j$. Each of those clusters sits on top of a dark
matter halo, and each halo gravitationally lenses the galaxies
behind it. $\obsij{R}$ is the model's prediction for the
\emph{average} lensing signal --- the excess surface mass density
--- that this bin of clusters would produce at projected distance
$R$ from their centres, if we could stack thousands of them and
average away the noise. It is built from three physical
ingredients, multiplied together and summed over every halo that
could possibly land in the bin:
\begin{itemize}
\item \textbf{How many halos are there?} The halo mass function
      $n(M,z)$ counts halos of mass $M$ at redshift $z$ per unit
      comoving volume; $\dVdz(z)\,\Omega(z)$ converts that into a
      count per unit redshift, weighted by how much sky area the
      survey covers at that redshift.
\item \textbf{Which of those halos actually end up in our bin?}
      Not every halo of mass $M$ at redshift $z$ is observed with
      richness in range $i$ and photo-$z$ in range $j$ --- richness
      and photo-$z$ are noisy proxies for the true $(M,z)$. The
      selection function $S_{ij}(\lnM,z)$ is exactly this
      probability, precomputed once by \texttt{sel\_function.py}
      and factorised there as $S_{ij}=S_i(\lnM,z)\,K_j(z)$: a
      richness-selection piece $S_i$ that does the heavy lifting in
      mass, times a redshift-window piece $K_j$ that carves out
      bin $j$ (\texttt{src/models/sel\_function\_t.hh}).
\item \textbf{How strongly does each surviving halo lens?} That is
      the actual physics term, $\DSnfw(R,\lnM)$: the excess surface
      density of an NFW halo of mass $M$, evaluated at the radius
      $R$ we care about. This is the \emph{only} place $R$ enters
      the whole calculation.
\end{itemize}
Putting the three pieces together, the code
(\texttt{NOperatorSelRadial<DSigma1hWeight>} in
\texttt{src/models/n\_operator\_sel\_t.hh} and
\texttt{src/modules/num\_counts\_sel/lensing\_weights.hh}) evaluates,
for every bin $(i,j)$ and every radius $R$,
\begin{align}
N_{ij}(R) &= \int\! dz \int\! d\lnM\;
n(M,z)\,\dVdz(z)\,\Omega(z)\, S_{ij}(\lnM,z)\,
\DSnfw(R,\lnM),
\notag\\
\obsij{R} &= \frac{N_{ij}(R)}{N_{ij}},
\label{eq:full}
\end{align}
where $N_{ij} = \int dz\int d\lnM\; n\,\dVdz\,\Omega\,S_{ij}$ is the
same bin's halo count, with no $R$-dependence at all (this is what
the \texttt{NumCounts} module computes). In words,
eq.~\eqref{eq:full} says: \emph{count the
halos that land in bin $(i,j)$, weight each one by how it lenses
at radius $R$, and average.}

\subsection{Why this is slow}

Nothing in eq.~\eqref{eq:full} looks unusual --- it is a completely
standard halo-model stacking integral. The expense is purely
computational: the code cannot do the $(\lnM,z)$ quadrature once
and reuse it, because the integrand at radius $R_1$ and the
integrand at radius $R_2$ look different (they differ in the
$\DSnfw(R,\lnM)$ factor), so \emph{the whole 2-D quadrature is
rerun from scratch at every point of the radius grid}, for every
bin. If there are, say, $20$ radii and $12$ richness$\times$redshift
bins, that is $240$ independent 2-D quadratures where, as we are
about to see, one would (almost) do.

\subsection{A free win: the halo does not know which redshift bin it is in}
\label{sec:zwin}

Here is the key observation, and it is worth stating carefully
because it is what makes \S\ref{sec:zwin} an \emph{exact}
simplification rather than an approximation: \emph{look again at
eq.~\eqref{eq:full} and ask which factors actually depend on $z$.}
$n(M,z)$, $\dVdz(z)$, $\Omega(z)$, and $S_{ij}(\lnM,z)$ all do, by
construction. But $\DSnfw(R,\lnM)$ --- the one factor that also
depends on $R$ --- does not carry a $z$-argument at all in this
model. That is because the NFW profile used here fixes the
concentration at a constant, $c=4$, and defines the halo boundary
with a fixed reference density $\rho_{\rm crit,0}$ (optionally
rescaled by a constant $\Omega_m$ factor, but never by
$\rho_{\rm crit}(z)$; see \texttt{src/models/nfw\_dsigma\_mis.hh}).
Physically: a cluster's lensing \emph{shape} in this model is
entirely a function of its mass, full stop --- redshift only ever
enters through \emph{how many} halos of that mass exist and how
likely they are to fall in bin $(i,j)$, never through what the
profile itself looks like.

Once you notice that, the reduction is just Fubini's theorem ---
swap the order of integration and do the $z$-integral first, since
$\DSnfw(R,\lnM)$ has nothing to do with $z$ and can be pulled
outside it:
\begin{equation}
\boxed{\;
\begin{aligned}
\widehat W_{ij}(\lnM) &\equiv \int\! dz\;
n(M,z)\,\dVdz(z)\,\Omega(z)\, S_{ij}(\lnM,z),\\[2pt]
N_{ij}(R) &= \int\! d\lnM\; \widehat W_{ij}(\lnM)\,\DSnfw(R,\lnM).
\end{aligned}
\;}
\label{eq:reduction}
\end{equation}
In words: build a mass-only weight $\widehat W_{ij}(\lnM)$ once per
bin --- how many halos of mass $M$, integrated over the whole
redshift window, end up in bin $(i,j)$ --- and then the radius grid
only ever touches a \emph{1-D} mass integral. This costs nothing
extra to implement (it reuses the same mass--redshift quadrature
nodes the code already visits) and is exact to machine precision,
because it is nothing more than doing the integral in a different
order. It should be implemented regardless of everything that
follows. What it does \emph{not} solve is the remaining
$R$-dependence sitting inside $\DSnfw(R,\lnM)$ itself --- that is
the real obstruction, and the rest of this note is about it.

\section{What is left: a stack of different-sized NFW profiles}
\label{sec:setup}

\subsection{Why \texorpdfstring{$R$}{R} still cannot be pulled out}

Every NFW profile has the same functional shape, but that shape is
stretched by a different amount for every halo mass: heavier
halos have a larger scale radius $\rs(M)$, so their $\DSnfw$ curve
is a stretched-out version of a lighter halo's curve, not a
rescaled copy of the \emph{same} curve at the same $x=R/\rs$.
Concretely,
\begin{equation}
\DSnfw(R,\lnM) = A(M)\, g\!\left(\frac{R}{\rs(M)}\right),
\qquad
g(x) = \text{Wright \& Brainerd (2000) NFW shape function},
\label{eq:nfw_shape}
\end{equation}
with $\rs(M) = r_{200}(M)/c$, $r_{200}(M) = [3M/(800\pi\rho_{\rm
crit})]^{1/3}$, $A(M) = 2\rs(M)\,\delta_c\,\rho_{\rm crit}$, and
$\delta_c = \tfrac{200}{3}c^3/[\ln(1+c)-c/(1+c)]\approx5.27\times10^3$
at $c=4$. If $\rs$ did not depend on $M$, eq.~\eqref{eq:reduction}
would already be done: $g(R/\rs)$ would come straight out of the
mass integral and $N_{ij}(R)$ would be a single evaluation of $g$
times a mass-integrated amplitude. It is precisely because
\emph{different halos in the same bin have different scale radii}
that the mass integral mixes together many differently-stretched
copies of the same curve, and the result is a genuinely new shape
--- broader and shallower than any individual NFW profile, the way
adding together several Gaussians of different widths gives you
something that is no longer a single Gaussian.

\subsection{A change of variable that makes the mixing explicit}

Since $\ln\rs(M) = \tfrac13\lnM + {\rm const}$, it is convenient to
stop thinking in $M$ and think directly in $y\equiv\ln\rs(M)$ ---
the log-scale-radius is what actually controls the profile shape.
Define the pure kernel
\begin{equation}
K(R,y) \equiv A\big(M(y)\big)\, g\!\left(R\,e^{-y}\right)
= C\,e^{y}\, g\!\left(R\,e^{-y}\right),
\qquad C \equiv 2\,\delta_c\,\rho_{\rm crit},
\label{eq:kernel}
\end{equation}
i.e.\ ``the lensing profile of a halo with log-scale-radius $y$,
evaluated at $R$.'' Relabelling the weight of eq.~\eqref{eq:reduction}
in the same variable, $\Wcal_{ij}(y)\,dy \equiv \widehat
W_{ij}(\lnM)\,d\lnM$ (a constant Jacobian, $d\lnM=3\,dy$), the
mass integral becomes
\begin{equation}
N_{ij}(R) = \int\! dy\; \Wcal_{ij}(y)\, K(R,y).
\label{eq:reduced_1d}
\end{equation}
This is exactly the same integral as eq.~\eqref{eq:reduction} ---
nothing has been approximated --- but written this way it is
transparent \emph{why} it is expensive: $N_{ij}(R)$ is a weighted
average of the kernel $K(R,y)$ over the distribution of
scale radii $\Wcal_{ij}(y)$ present in bin $(i,j)$, and that
average has to be redone at every $R$ because $K(R,y)$'s shape in
$y$ changes with $R$.

\subsection{How wide is the mixing, in practice?}

The three approximations below all work by exploiting that
$\Wcal_{ij}(y)$ is \emph{narrow}: most of the halos contributing
to a given richness$\times$redshift bin have a fairly similar
mass, hence a fairly similar $\rs$. Quantitatively, a bin with mass
scatter $\sigma_{\lnM}\sim0.3$--$0.6$ (combining intrinsic
mass--richness scatter and the finite width of the bin itself)
gives $\sigma_y=\sigma_{\lnM}/3\sim0.1$--$0.2$ in log-scale-radius
--- a $10$--$20\%$ spread in $\rs$, not a factor of a few. That
narrowness is exactly what makes it plausible to replace the full
mixture with a compact approximation, and it sets the size of the
errors quoted for each approach below.

\section{Three ways to pull \texorpdfstring{$R$}{R} out of the mass integral}
\label{sec:approaches}

The three ideas below go from simplest-and-cheapest to
most-robust-and-general. Each one answers the same question ---
``can I replace the weighted average over $y$ in
eq.~\eqref{eq:reduced_1d} with something that does not need to be
recomputed at every $R$?'' --- with a different strategy.

\subsection{Idea 1: pretend the whole bin is one effective halo}
\label{sec:approach1}

\paragraph{The idea.} If every halo in bin $(i,j)$ had exactly the
same scale radius, $N_{ij}(R)$ would just be (halo count) times
(the one shape function). Bin $(i,j)$'s halos do not all have the
same scale radius, but they are close (\S\ref{sec:setup}), so the
simplest thing to try is to replace the whole distribution by the
single halo of \emph{mean} log-scale-radius --- equivalently, the
halo of geometric-mean mass --- present in the bin:
\begin{equation}
\boxed{\;
N_{ij}(R) \;\approx\; \Aeff_{ij}\; g\!\left(\frac{R}{\reffi{ij}}\right),
\qquad
\yeff_{ij} = \frac{\int dy\,\Wcal_{ij}(y)\, y}{\int dy\,\Wcal_{ij}(y)},
\qquad
\reffi{ij} = e^{\yeff_{ij}},
\qquad
\Aeff_{ij} = C\, e^{\yeff_{ij}}.
\;}
\label{eq:approach1}
\end{equation}
Two things are worth being precise about, because both are easy to
get wrong. First, $\yeff_{ij}$ is the \emph{plain} mean of $\ln\rs$
under the bin's own mass weight $\Wcal_{ij}$ --- \emph{not} weighted
by the amplitude $Ce^y$. Second, $\Aeff_{ij}$ is the amplitude
\emph{evaluated at} that mean scale radius, $Ce^{\yeff_{ij}}$ ---
\emph{not} the amplitude averaged separately over the bin,
$\int dy\,\Wcal_{ij}(y)\,Ce^y$ (these differ, because $e^y$ is
convex: the average of $e^y$ is always $\ge e^{\text{average of
}y}$, by Jensen's inequality). This particular pairing --- plain
mean of $y$, amplitude evaluated there --- is not a stylistic
choice: it is exactly the pairing for which the kernel identity of
eq.~\eqref{eq:deriv_identity} below makes the first correction term
of Idea 2 (\S\ref{sec:approach2}) vanish identically, so Idea 1
\emph{is} Idea 2 truncated at the very first term, and any other
pairing of ``an amplitude'' with ``a scale radius'' would not have
that property. If a physically-labelled effective mass is wanted,
$M_{\rm eff}=\exp(\langle\lnM\rangle_{\Wcal_{ij}})$ --- the
geometric-mean halo mass of the bin --- reproduces $\reffi{ij}$ and
$\Aeff_{ij}$ exactly, since $\ln\rs(M)$ is linear in $\lnM$; the
concentration stays fixed at $c=4$, no $c_{\rm eff}$ needed.

\paragraph{Cost.} Two numbers per bin, computed once; one
evaluation of $g$ per radius --- as cheap as it gets.

\paragraph{Where the error comes from.} Squeezing a whole
distribution of stretched NFW curves into one curve necessarily
loses the fact that the real stack is a bit broader/shallower than
any single member of it. That broadening is the leading term Idea
1 throws away (formally, the $n=2$ term of Idea 2's series,
\S\ref{sec:approach2}): roughly $\sigma_y^2\times\mathcal{O}(1$--$5)$
in size. \S\ref{sec:validation} measures this directly: for a
representative mass weight, Idea 1 is already off by
$\sim3\%$ at $\sigma_{\lnM}=0.3$ and grows to $\sim12\%$ at
$\sigma_{\lnM}=0.6$; the error is smallest at small $R$ and grows
(roughly monotonically, then saturating) toward the outer part of
the radial range, rather than peaking at any particular $R$ ---
this is the approach to avoid if $1\%$-level accuracy matters
anywhere in the pipeline.

\paragraph{When it breaks.} If $\Wcal_{ij}(y)$ has two separate
humps or a heavy tail (e.g.\ a bin with unusually large
mass--richness scatter, or contamination from the low-richness
Poisson tail), no single curve can represent it well anywhere, and
the good agreement near $R\approx\reffi{ij}$ does not extend to
the edges of the radial range.

\subsection{Idea 2: keep the correction terms Idea 1 threw away (recommended first cut)}
\label{sec:approach2}

\paragraph{The idea.} Idea 1 is the zeroth term of a systematic
expansion; rather than stopping there, expand the kernel
eq.~\eqref{eq:kernel} in a Taylor series around $y=\yeff_{ij}$ and
keep going. The key algebraic fact that makes this work cleanly is
that differentiating the kernel with respect to $y$ is the same
operation as differentiating the shape function $g$ with respect
to $\ln x$: writing $\Lop\equiv\partial/\partial\ln x$,
\begin{equation}
\frac{\partial^n}{\partial y^n}\Big[C e^y g(Re^{-y})\Big]
= C e^y\,(1-\Lop)^n g(x),
\qquad x \equiv R e^{-y}.
\label{eq:deriv_identity}
\end{equation}
Because of this identity, Taylor-expanding in $y$ is \emph{exact,
order by order} --- nothing is dropped except the terms you
explicitly choose to truncate:
\begin{equation}
\boxed{\;
N_{ij}(R) = \Aeff_{ij} \sum_{n\ge0} \frac{\mu_n^{(ij)}}{n!}\,
\big[(1-\Lop)^n g\big](\hat x_{ij}),
\qquad
\hat x_{ij} = \frac{R}{\reffi{ij}},
\qquad
\mu_n^{(ij)} = \frac{\int dy\,\Wcal_{ij}(y)\,(y-\yeff_{ij})^n}
{\int dy\,\Wcal_{ij}(y)},
\;}
\label{eq:approach2}
\end{equation}
with $\Aeff_{ij}$ and $\yeff_{ij}$ exactly as defined in
eq.~\eqref{eq:approach1}: $\mu_n^{(ij)}$ is the \emph{plain}
$n$-th central moment of $\ln\rs$ under the bin's mass weight ---
there is no amplitude factor inside this integral, unlike
$\Aeff_{ij}$ itself. This is the pairing that makes
$\mu_0^{(ij)}=1$ and $\mu_1^{(ij)}=0$ automatic (the latter is
exactly why $\yeff_{ij}$ was defined as a plain mean in
eq.~\ref{eq:approach1} --- an amplitude-weighted mean would leave a
nonzero, uncancelled linear term here). The $n=0$ term is Idea 1; every higher $n$
is a correction that captures more of the true spread of scale
radii in the bin. The needed operators are closed-form
combinations of the NFW shape derivatives ($g'\equiv dg/dx$):
\begin{equation}
(1-\Lop)^0 g = g,
\qquad
(1-\Lop)^2 g = g - x g' + x^2 g'',
\qquad
(1-\Lop)^3 g = g - x g' - x^3 g'''.
\label{eq:operators}
\end{equation}

\paragraph{Cost.} $3$--$4$ numbers per bin
($\Aeff_{ij},\reffi{ij},\mu_2^{(ij)}$, optionally $\mu_3^{(ij)}$),
accumulated in the \emph{same} mass-quadrature pass already needed
for eq.~\eqref{eq:reduction} --- computing a moment costs no more
than computing the weight itself. Per radius: evaluate
$g,g',g''$(,$g'''$) at $\hat x_{ij}$ and add up 3--4 terms;
$\hat x_{ij}$ is different for every bin, so the total work is
$N_{\rm radii}\times N_{\rm bins}$ cheap analytic evaluations, with
\emph{zero} additional quadrature.

\paragraph{How accurate is truncating here?} Stopping at $n=2$
leaves out $n=4$ as the next contribution (the $n=3$ term is
already close to zero for a roughly symmetric $\Wcal_{ij}$); for a
Gaussian-like weight, $\mu_4\approx3\mu_2^2$, so the relative error
left behind is of order $(\mu_2^{(ij)})^2$, i.e.\ two more powers of
$\sigma_y^2$ down from Idea 1's already-small error.
\S\ref{sec:validation} confirms this is dramatic in practice: for
the same representative weight where Idea 1 was off by $3$--$12\%$,
truncating at $n\le2$ brings the error down to
$0.03$--$0.6\%$ across $\sigma_{\lnM}=0.3$--$0.6$. \textbf{In
short: $n\le2$ is already the workhorse choice.} Adding
$\mu_3^{(ij)}$ ($n=3$) did \emph{not} help further in that test ---
consistent with $\mu_3$ being small and slightly negative for a
near-symmetric weight, so the $n=3$ term is comparable in size to
noise from the already-tiny $n\ge4$ remainder rather than a
systematic improvement; keep $n\le2$ unless a specific bin is
visibly skewed.

\paragraph{When it breaks.} This is a series in powers of
$\sigma_y^2$, so it only converges quickly while $\Wcal_{ij}(y)$ is
reasonably compact and unimodal --- the same condition Idea 1
needed, just enforced to higher order. A strongly bimodal or
heavy-tailed weight inflates every moment $\mu_3,\mu_4,\dots$
comparably, and adding more terms stops helping. As with Idea 1,
the residual grows (slowly, and saturating) toward the outer part
of the radial range rather than peaking at one particular $R$
(Fig.~\ref{fig:validation_profiles}).

\subsection{Idea 3: expand the halo shape itself, not the bin's weight}
\label{sec:approach3}

\paragraph{The idea.} Ideas 1 and 2 both expand \emph{the bin's
distribution of scale radii} and therefore only work well when
that distribution is narrow and well-behaved. Idea 3 flips the
strategy: instead of approximating $\Wcal_{ij}(y)$, approximate the
kernel $K(R,y)$ itself --- a fixed, bin-\emph{independent} function
--- as a short sum of basis functions in $y$. Because $K(R,y)$ is a
smooth, analytic function of $y$ for any fixed $R$, it is very well
approximated on the full range of scale radii spanned by
\emph{all} the bins at once, $y\in[y_{\min},y_{\max}]$ (rescaled
affinely to $\tilde y\in[-1,1]$), by a short Chebyshev series:
\begin{equation}
\boxed{\;
K(R,y) \approx \sum_{k=0}^{K} u_k(R)\, T_k(\tilde y),
\qquad
N_{ij}(R) \approx \sum_{k=0}^{K} c_k^{(ij)}\, u_k(R),
\qquad
c_k^{(ij)} = \int\! dy\,\Wcal_{ij}(y)\, T_k(\tilde y).
\;}
\label{eq:approach3}
\end{equation}
The coefficients $u_k(R)$ depend only on the kernel and the
radius, not on the bin, so they are computed \emph{once, for every
bin simultaneously}. Each bin then just needs $K{+}1$ numbers
$c_k^{(ij)}$ --- projections of its own weight $\Wcal_{ij}$ onto the
same fixed basis functions --- accumulated in the same
mass-quadrature pass as Idea 2.

\paragraph{Cost.} $K{+}1\approx6$--$10$ numbers per bin (a bit more
bookkeeping than Idea 2's 3--4); $u_k(R)$ is shared across every
bin; per radius, a length-$(K{+}1)$ dot product.

\paragraph{How accurate is this?} Chebyshev coefficients of a
smooth, analytic function decay geometrically, so $K\approx6$--$10$
terms over a $y$-window covering roughly $6$ dex in $\ln M$
already gives $\sim10^{-3}$ accuracy, uniformly in $R$ ---
\emph{regardless of how wide, skewed, or multimodal $\Wcal_{ij}$
is}, because the approximation error lives entirely in the kernel,
never in the weight.

\paragraph{When it breaks.} Only by extrapolation --- if a bin's
mass range falls outside $[y_{\min},y_{\max}]$, or a radius falls
outside the tabulated range, the approximation degrades sharply.
Both are avoided simply by building the box from the union of
every bin's actual quadrature range. Otherwise this approach is
robust to exactly the pathologies (bimodal or heavy-tailed
$\Wcal_{ij}$) that break Ideas 1 and 2.

\section{Comparison and recommendation}
\label{sec:comparison}

\begin{center}
\small
\resizebox{\textwidth}{!}{%
\begin{tabular}{l c c c c}
\toprule
Idea & precompute/bin & per-radius cost & measured error (\S\ref{sec:validation}) & robust to skew/bimodal $\Wcal_{ij}$ \\
\midrule
1. One effective halo $(\Aeff,\reff)$ & 2 numbers & 1 $g$-eval & $3$--$12\%$ & no \\
2. Moment expansion ($n\le2$) & 3 numbers & 3 analytic terms & $0.03$--$0.6\%$ & no (asymptotic in $\sigma_y$) \\
3. Chebyshev kernel ($K=8$) & 9 numbers & 9-term dot product & $0.001$--$0.003\%$ & yes \\
\bottomrule
\end{tabular}%
}
\end{center}

Idea 2 with $n\le2$ is the one to build first: it slots directly
into the mass-quadrature loop that already has to compute
$\widehat W_{ij}(\lnM)$ (eq.~\ref{eq:reduction}) --- the only
change is accumulating one extra weighted moment alongside it ---
the NFW shape derivatives $g',g''$ are closed-form, and
\S\ref{sec:validation} shows it already buys three orders of
magnitude over Idea 1 for a representative weight. Idea 3 is the
more future-proof production choice, at roughly three times the
per-bin bookkeeping, and is the right fallback if the real
mass--richness relation $P(\lambda\mid M,z)$ turns out to be
skewed or multimodal for any bin (\S\ref{sec:validation} only
tests unimodal weights). Idea 1 is the cheapest of the three but
carries an irreducible few-percent-to-double-digit shape bias for
realistic bin widths; treat it as a sanity check or a starting
guess for fitting the other two, not as the production answer.

\section{Numerical validation}
\label{sec:validation}

\paragraph{What is tested here, and what is not yet.} Two separate
claims need checking: (i) is the NFW shape function $g(x)$ itself
implemented correctly, and (ii) given a correct $g(x)$, do the three
ideas actually reproduce the exact mass integral at the error levels
argued in \S\ref{sec:approaches}? Both are checked below. What is
\emph{not} checked yet is whether the real richness--mass relation
$P(\lambda\mid M,z)$ used by this pipeline's \texttt{sel\_function.py}
produces a mass weight $\Wcal_{ij}(y)$ that looks like the
representative one used here --- doing that requires reading the
actual selection-function tables out of a running pipeline sample,
which is a natural next step but a separate piece of work from the
math check below.

\paragraph{Setup.} All constants ($c=4$, $\rho_{\rm crit}=2.775\times
10^{11}\,\hMsun\,\hMpc^{-3}$) match
\texttt{src/models/nfw\_dsigma\_mis.hh} exactly. The mass weight
$\Wcal_{ij}(\lnM)$ is built as a Tinker-like halo-mass-function shape,
$(M/M_\star)^{-\alpha}e^{-(M/M_\star)^\beta}$ with $\alpha=0.9$,
$\beta=1.3$, $M_\star=2\times10^{14}\,\hMsun$, times a lognormal
mass--richness kernel of width $\sigma_{\lnM}$ centred on $M_0$; three
bins span the range quoted in \S\ref{sec:setup}: ``narrow''
($\sigma_{\lnM}=0.30$), ``medium'' ($\sigma_{\lnM}=0.45$), and
``wide'' ($\sigma_{\lnM}=0.60$). This is a stand-in for the true
$\Wcal_{ij}$, not a read-out of it (see above) --- but it is built
from the same physical ingredients (a steeply falling mass function
times a lognormal selection kernel) and spans the same
$\sigma_{\lnM}$ range the note has been arguing about throughout, so
it exercises exactly the mechanism the three ideas depend on. Script:
\texttt{docs/figs/make\_shear1h\_validation\_figs.py}.

\paragraph{Check (i): the shape function.} $g(x)$ was implemented
from the closed-form Wright \& Brainerd (2000) expression and
independently cross-checked against a from-scratch numerical
line-of-sight projection of the 3-D NFW density,
$\Sigma(R)=2\int_0^\infty\!\rho(\sqrt{R^2+z^2})\,dz$ and
$\Delta\Sigma(R)=\bar\Sigma(<R)-\Sigma(R)$, computed by direct
quadrature with no reference to the closed form at all. The ratio of
the two is constant to $1.6\times10^{-9}$ across
$x\in[0.05,10]$, including on both sides of the removable
singularity at $x=1$ --- the closed-form shape function is correct.

\paragraph{Check (ii): the three ideas against the exact mass
integral.} Table~\ref{tab:validation} and
Fig.~\ref{fig:validation_profiles} compare $N_{ij}(R)$, computed by
direct quadrature over $\lnM$ (eq.~\ref{eq:reduced_1d}), against
Ideas 1, 2 ($n\le2$ and $n\le3$), and 3 (Chebyshev order $K=8$) on a
$40$-point log-spaced radial grid, $R\in[0.05,20]\,\hMpc$.

\begin{table}[!htbp]
\centering
\small
\begin{tabular}{l c c c c c}
\toprule
Bin & $\sigma_{\lnM}$ & \multicolumn{2}{c}{Idea 1} & \multicolumn{1}{c}{Idea 2 ($n\le2$)} & Idea 3 ($K=8$) \\
& & max \% & rms \% & max \% & max \% \\
\midrule
narrow & 0.30 & 2.94 & 1.89 & 0.028 & 0.002 \\
medium & 0.45 & 6.94 & 4.66 & 0.179 & 0.003 \\
wide   & 0.60 & 12.29 & 8.51 & 0.611 & 0.001 \\
\bottomrule
\end{tabular}
\caption{Fractional residual of each idea against the exact mass
integral, over the full radial grid, for the representative weight of
\S\ref{sec:validation}. Idea 2's $n\le3$ column is omitted because it
did not improve on $n\le2$ here (\S\ref{sec:approach2}); Idea 3's rms
is $\le0.001\%$ in every bin.}
\label{tab:validation}
\end{table}

\begin{figure}[!htbp]
\centering
\includegraphics[width=\textwidth]{figs/shear1h_validation_profiles.png}
\caption{Top: exact $N_{ij}(R)$ (black) against Ideas 1, 2 ($n\le3$),
3, for the narrow/medium/wide representative bins. Bottom: fractional
residual. Idea 3 is essentially exact at this scale; Idea 1's error
grows visibly with bin width and grows monotonically (then
saturates) toward large $R$, rather than peaking at $R\approx\reffi{ij}$
as originally guessed in \S\ref{sec:approach1} --- corrected here now
that it has actually been measured.}
\label{fig:validation_profiles}
\end{figure}

\begin{figure}[!htbp]
\centering
\includegraphics[width=0.62\textwidth]{figs/shear1h_validation_scaling.png}
\caption{Maximum fractional residual vs.\ bin width $\sigma_{\lnM}$,
one curve per idea. Idea 1's error grows roughly as $\sigma_{\lnM}^2$,
consistent with it being the leading truncated term of Idea 2's
series (\S\ref{sec:approach1}); Idea 2 improves on it by
$\sim2$ orders of magnitude across the same range; Idea 3 is flat and
nearly a further order of magnitude below Idea 2.}
\label{fig:validation_scaling}
\end{figure}

The results confirm the qualitative argument of
\S\ref{sec:approaches} and sharpen it quantitatively: Idea 1 is not
safe to use directly once $\sigma_{\lnM}\gtrsim0.4$ (already
$\sim7\%$ off at $\sigma_{\lnM}=0.45$), Idea 2 at $n\le2$ removes
essentially all of that error (sub-percent even at the widest bin
tested), and Idea 3 is indistinguishable from exact at the
resolution of this test. One correction to \S\ref{sec:approach2}'s
original estimate: an earlier draft of this note expected Idea 2 to
need $n\le3$ for the widest bins; in practice $n\le2$ already
outperforms that expectation by an order of magnitude, and adding
$n=3$ made things very slightly worse rather than better (consistent
with $\mu_3$ being small and sign-uncertain for a near-symmetric
weight, so it mostly adds noise once $\mu_2$ has already done the
work).

\paragraph{Remaining checklist before production use.} Items 1 and 3
below are now done with the real pipeline in
\S\ref{sec:real_validation}; item 2 is the one piece still open.
\begin{enumerate}
\item \textbf{[Done]} Read $\Wcal_{ij}(\lnM)$ out of an actual
      pipeline sample: \S\ref{sec:real_validation} pulls it directly
      from a running \texttt{cosmosis} sample (real HMF, distances,
      $S_{ij}$), cross-checked against \texttt{NumCountsSel.so}
      to $<0.1\%$.
\item Confirm eq.~\eqref{eq:reduction} reproduces the current
      \texttt{DSigma1hSel} output bit-for-bit for a fixed bin and
      radius grid (it must, to machine precision, since it is a
      reordering of the integral, not an approximation) before
      swapping in Idea 2 on top of it --- still open; what
      \S\ref{sec:real_validation} confirms is the zeroth moment
      ($\int d\lnM\,\widehat W_{ij}$ vs. \texttt{NumCountsSel.so}),
      not yet the full $R$-dependent \texttt{DSigma1hSel.so} output
      point by point.
\item \textbf{[Done]} Re-run the comparison with the real
      $\Wcal_{ij}$: \S\ref{sec:real_validation}, together with the
      miscentered mixture and a cross-check against
      \texttt{Shear1hMisSel.so}.
\end{enumerate}

\section{Extending to the miscentered profile}
\label{sec:mis}

Everything above treats the pure centred term,
\texttt{DSigma1hSel}/\texttt{DSigma1hWeight}. The production
pipeline's actual lensing observable mixes in a miscentered
component (\texttt{DSigma1hMisSel},
\texttt{src/modules/num\_counts\_sel/lensing\_weights.hh}):
\begin{equation}
\Delta\Sigma_{\rm cl}(R,M) = (1-f_{\rm mis})\,\DSnfw(R,M)
+ f_{\rm mis}\,\Delta\Sigma_{\rm mis}\big(R,M;\,\tau_{\rm mis}R_\lambda(\lambda^{\rm ob})\big),
\label{eq:mixture}
\end{equation}
with DES-Y3 fiducial $f_{\rm mis}=0.22$, $\tau_{\rm mis}=0.17$, and
$R_\lambda(\lambda^{\rm ob})=(\lambda^{\rm ob}/100)^{0.2}$. A
fraction $f_{\rm mis}$ of clusters are offset from the true halo
centre; $\Delta\Sigma_{\rm mis}$ is the offset-averaged profile,
read from the precomputed table
\texttt{data/nfw\_off\_center/*\_gamma\_*.txt}
(\texttt{NFW\_DSIGMA\_MIS}, the ``gamma'' Rayleigh-type kernel of
Kelly et al.\ 2023) rather than a closed form.

\subsection{The good news: the same machinery applies, with one simplification}

Two facts make this an easy extension rather than a new derivation.
First, $R_\lambda$ depends only on the richness bin, not on mass or
redshift, so it is a \emph{fixed constant for a given bin} --- call
it $\rmis \equiv \tau_{\rm mis}R_\lambda$. Second, and more usefully:
\S\ref{sec:approach2}'s Taylor expansion never actually needed the
amplitude/shape split $A(M)g(x)$ --- it only needed \emph{ordinary
derivatives with respect to $y$} of whatever profile sits inside the
mass integral. Define, for a fixed radius $R$ and bin-level $\rmis$,
\begin{equation}
\Phi(R,y) \equiv \Delta\Sigma_{\rm cl}\big(R,\,M(y)\big)
= (1-f_{\rm mis})\,\DSnfw(R,M(y)) + f_{\rm mis}\,\Delta\Sigma_{\rm mis}(R,M(y);\rmis),
\label{eq:Phi}
\end{equation}
and Idea 2 becomes, with no change in logic at all:
\begin{equation}
\boxed{\;
N_{ij}(R) \approx \sum_{n\ge0} \frac{\mu_n^{(ij)}}{n!}\,\Phi^{(n)}(R,\yeff_{ij}),
\qquad
\Phi^{(n)} \equiv \frac{\partial^n\Phi}{\partial y^n},
\;}
\label{eq:approach2_general}
\end{equation}
with $\yeff_{ij}$ and $\mu_n^{(ij)}$ \emph{exactly as defined before}
(eqs.~\ref{eq:approach1}--\ref{eq:approach2}) --- they are moments of
the mass weight $\Wcal_{ij}$ alone and know nothing about
miscentering. $\Phi^{(n)}(R,\yeff_{ij})$ is computed numerically:
evaluate $\Phi(R,y)$ on a small grid of $y$ around $\yeff_{ij}$ (a
handful of table look-ups) and read off derivatives from a local
spline. Idea 1 is $n=0$; Idea 2 keeps $n=2$. eq.~\eqref{eq:deriv_identity}
is the special case of eq.~\eqref{eq:approach2_general} where $\Phi$
happens to factor as $Ce^yg(x)$ and the derivatives happen to have a
closed form --- eq.~\eqref{eq:approach2_general} is the more general
statement, and it is what \S\ref{sec:real_validation} actually
implements, directly against the real tables, with \emph{no}
amplitude/shape factoring anywhere. Idea 3 carries over the same
way: Chebyshev-fit $\Phi(R,y)$ itself (not a factored kernel) in
$y$; the one change is that the fit is no longer shared globally
across richness bins, since $\rmis$ (and hence $\Phi$) now differs
bin to bin, so $u_k(R)$ must be refit per bin (still shared across
redshift bins at fixed richness, since $\rmis$ does not depend on
$z$).

\subsection{A validation subtlety: two independently-sourced tables}
\label{sec:mis_subtlety}

$\DSnfw(R,M)$ in production comes from a Python-computed table
(\texttt{haloModel/dSigma\_nfw}, written by
\texttt{halo\_model\_cosmosis.py}); $\Delta\Sigma_{\rm mis}$ comes
from the independent C++ table read by \texttt{NFW\_DSIGMA\_MIS}.
Confirming eq.~\eqref{eq:mixture} reduces correctly to the centred
profile as the offset becomes negligible requires knowing both
tables sit on the same absolute physical scale --- which the
production code presumably guarantees, but which this note has not
independently re-derived from \texttt{halo\_model\_cosmosis.py}. For
\S\ref{sec:real_validation}'s validation, this is sidestepped by
defining
\begin{equation}
\Delta\Sigma_{\rm mis}(R,M;\rmis) \equiv \DSnfw^{\rm real}(R,M)\times
{\rm ratio}(x,x_{\rm mis}),
\qquad
{\rm ratio}(x,x_{\rm mis}) \equiv \frac{g_{\rm mis}(x,x_{\rm mis})}{g(x)},
\label{eq:mis_ratio_def}
\end{equation}
with $x=R/\rs(M)$, $x_{\rm mis}=\rmis/\rs(M)$. By construction
${\rm ratio}\to1$ for $x\gg x_{\rm mis}$ (verified against the real
gamma table in \S\ref{sec:real_validation}): the real centred table,
rescaled by this dimensionless, table-derived suppression factor.
This guarantees the correct
centred limit regardless of any independent normalisation
convention difference between the two source tables, at the cost of
not being a literal re-implementation of \texttt{NFW\_DSIGMA\_MIS}'s
own internal amplitude. \S\ref{sec:real_validation} quantifies the
resulting $\sim4\%$ shape difference from the real
\texttt{Shear1hMisSel.so} output --- small enough that it does not
change which idea to recommend, but real enough to flag rather than
hide.

\section{Validation against the real production pipeline}
\label{sec:real_validation}

Rather than a synthetic stand-in, this section pulls every ingredient
from an actual running \texttt{cosmosis} sample: the real halo mass
function, the real comoving-distance table, the real
richness-selection function $S_{ij}(\lnM,z)$
(\texttt{sel\_function.py}), and the real
\texttt{haloModel/dSigma\_nfw} table --- plus the real
miscentering table for $\Delta\Sigma_{\rm mis}$
(\S\ref{sec:mis_subtlety}). Script:
\texttt{docs/figs/real\_pipeline\_extract.ini} +
\texttt{docs/figs/make\_real\_pipeline\_validation\_figs.py}; the
pipeline is
\texttt{consistency} $\to$ \texttt{GrowthFactor} $\to$ \texttt{cp\_camb}
$\to$ \texttt{MfTinker} $\to$ \texttt{halo\_model} $\to$
\texttt{average\_sigma\_crit\_inv} $\to$ \texttt{sel\_function} $\to$
\texttt{NumCountsSel} $\to$ \texttt{Shear1hMisSel},
run with the \texttt{test} sampler at the fiducial widePlanck point
from \texttt{des-cluster-nersc/cosmosis-models/}, restricted to the
four richness bins in the $z\in[0.20,0.35]$ block
($\lambda^{\rm ob}$ centres $25,\,37.5,\,52.5,\,130$) where the
mixture weight's default \texttt{lob\_centers} maps onto
\texttt{bin\_index} unambiguously.\footnote{\texttt{DSigma1hMisWeight}'s
\texttt{lob\_centers\_} defaults to length 4 (the DES-Y3 richness
bins), but \texttt{sel\_function}'s \texttt{bin\_index} runs 0--11
(4 richness $\times$ 3 redshift blocks); for \texttt{bin\_index}
$\ge4$ \texttt{set\_bin} is a no-op unless the ini overrides
\texttt{lob\_centers} to length 12. This is an orthogonal,
out-of-scope observation about the production ini, not something
this note's factorisation approach depends on --- avoided here by
simply validating on the unambiguous bins.}

\paragraph{Check (i): does this note's $\Wcal_{ij}$ match the real pipeline?}
A direct Python re-implementation of $n(M,z)$, $\dVdz(z)$,
$\Omega(z)$, and $S_{ij}(\lnM,z)$ against the sampled datablock,
integrated over $z$ \emph{and} $\lnM$ with fixed Gauss-Legendre
quadrature (64 and 96 nodes respectively --- this project's own
integrator convention, not a plain trapezoid grid) via
eq.~\eqref{eq:reduction}, is compared against the real
\texttt{NumCountsSel.so} output for each bin:
\begin{center}
\small
\begin{tabular}{c c c c}
\toprule
bin & $\int d\lnM\,\widehat W_{ij}$ (Python, fixed-GL) & \texttt{NumCountsSel.so} & ratio \\
\midrule
0 & $1.3804\times10^{3}$ & $1.3807\times10^{3}$ & $0.9998$ \\
1 & $5.1386\times10^{2}$ & $5.1383\times10^{2}$ & $1.0000$ \\
2 & $1.4148\times10^{2}$ & $1.4147\times10^{2}$ & $1.0001$ \\
3 & $9.1669\times10^{1}$ & $9.1733\times10^{1}$ & $0.9993$ \\
\bottomrule
\end{tabular}
\end{center}
Agreement to $<0.1\%$, with residuals now split roughly evenly above
and below $1.0$ (rather than the persistent same-sign bias a coarser
trapezoid grid would leave) --- consistent with the residual being
\texttt{NumCountsSel.so}'s own \texttt{eps\_rel}=$1.5\times10^{-3}$
Cuhre tolerance, not a resolution limit on the Python side. The exact
reduction of
eq.~\eqref{eq:reduction} and this note's replica of every ingredient
in it are both confirmed directly against the production code.

\paragraph{Check (ii): the mixture against \texttt{Shear1hMisSel.so}.}
Using $\Wcal_{ij}$ from check (i) and eq.~\eqref{eq:mixture} (with
$\Delta\Sigma_{\rm mis}$ per eq.~\ref{eq:mis_ratio_def}), this note's
own exact quadrature (eq.~\ref{eq:reduced_1d}) is compared, in
\emph{shape} (each curve normalised to its own $R=0.2\,\hMpc$ point,
since \S\ref{sec:mis_subtlety} already disclosed the two tables'
amplitudes are not independently re-derived), against the real
\texttt{Shear1hMisSel.so} output: maximum shape deviation
$3.6$--$4.2\%$ across the four bins. Small, and consistent with the
disclosed modelling choice in eq.~\eqref{eq:mis_ratio_def} rather
than with an error in the radial-factorisation approach itself ---
which is what the next comparison isolates.

\paragraph{Ideas 1, 2, and 3 against this note's own exact quadrature.}
Table~\ref{tab:real_validation} and
Fig.~\ref{fig:real_validation_profiles} compare the full mixture
$N_{ij}(R)$ (eq.~\ref{eq:approach2_general}) against all three
ideas, every one evaluated directly off the real tables with no
amplitude/shape factoring, on the real \texttt{Shear1hMisSel} radial
grid ($R\in[0.2,5]\,\hMpc$, 10 points). Idea 3 uses $K=8$ and refits
$u_k(R)$ per richness bin (\S\ref{sec:mis}), since $\rmis$ now
differs bin to bin.

\begin{table}[!htbp]
\centering
\small
\resizebox{\textwidth}{!}{%
\begin{tabular}{c c c c c c c c}
\toprule
bin & $\lambda^{\rm ob}$ & \multicolumn{2}{c}{Idea 1} & \multicolumn{2}{c}{Idea 2 ($n\le2$)} & \multicolumn{2}{c}{Idea 3 ($K=8$)} \\
& & max \% & rms \% & max \% & rms \% & max \% & rms \% \\
\midrule
0 & 25  & 9.75 & 6.53 & 0.365 & 0.238 & 0.001 & 0.001 \\
1 & 38  & 6.15 & 3.94 & 0.288 & 0.134 & 0.004 & 0.003 \\
2 & 52  & 3.91 & 2.43 & 0.093 & 0.050 & 0.002 & 0.001 \\
3 & 130 & 4.91 & 2.95 & 0.149 & 0.072 & 0.002 & 0.001 \\
\bottomrule
\end{tabular}%
}
\caption{Fractional residual of Ideas 1--3 against this note's own
exact quadrature of the real-table mixture, for the four real,
unambiguous richness bins of \S\ref{sec:real_validation}. Idea 3 is
$2$--$3$ orders of magnitude more accurate than Idea 2, and $3$--$4$
orders of magnitude more accurate than Idea 1, at real-pipeline
resolution --- it is the one to build if the extra per-bin
bookkeeping (\S\ref{sec:approach3}) is not a concern.}
\label{tab:real_validation}
\end{table}

\begin{figure}[!htbp]
\centering
\includegraphics[width=\textwidth]{figs/real_validation_profiles.png}
\caption{Real production tables: exact mixture (black) vs.\ Idea 1,
Idea 2 ($n\le2$), and Idea 3 ($K=8$), for the four real, unambiguous
richness bins. Bottom: fractional residual. Idea 3 is
indistinguishable from exact at this scale; Idea 2 sits at or below
$0.4\%$; Idea 1 reaches $5$--$10\%$ at the largest radii.}
\label{fig:real_validation_profiles}
\end{figure}

\begin{figure}[!htbp]
\centering
\includegraphics[width=0.62\textwidth]{figs/real_validation_scaling.png}
\caption{Maximum fractional residual vs.\ richness-bin centre, real
pipeline. Idea 3 sits $2$--$4$ orders of magnitude below Idea 1 and
consistently below Idea 2 across the full richness range tested,
matching the synthetic-weight result of \S\ref{sec:validation}.}
\label{fig:real_validation_scaling}
\end{figure}

This closes the loop opened in \S\ref{sec:validation}: the same
qualitative and quantitative conclusion --- Idea 3, then Idea 2, are
successively more accurate than Idea 1 by orders of magnitude, at
modest extra implementation cost --- holds with the real mass
function, the real selection function, and the real halo-model
table, for the full miscentered mixture, not just the synthetic
centred-only illustration. The remaining open item is
eq.~\eqref{eq:mis_ratio_def}'s $\sim4\%$ shape offset from the true
production mixture, which is a modelling choice made explicit in
\S\ref{sec:mis_subtlety} for lack of the Python
\texttt{dSigma\_nfw} table's internal amplitude convention, not a
property of the radial factorisation itself.

\section{What this buys the pipeline}
\label{sec:timing}

The point of all this is wall-clock time, so it is worth being
concrete about the current budget. Running the full production
pipeline (\texttt{mock\_mcmc\_cp\_camb.ini}, all 12 richness/redshift
bins, \texttt{test} sampler, one evaluation at the fiducial
widePlanck point) gives:
\begin{center}
\small
\begin{tabular}{l r}
\toprule
module & time [s] \\
\midrule
\texttt{MfTinker}                  & 0.142 \\
\texttt{halo\_model}               & 0.140 \\
\texttt{sel\_function}             & 0.191 \\
\texttt{NumCountsSel}              & 0.095 \\
\texttt{Shear1hMisSel}             & \textbf{0.458} \\
\texttt{b\_sel\_marg}              & 0.067 \\
\texttt{bsel}                      & 0.019 \\
\texttt{shear\_prj\_frozen\_physics} & 0.081 \\
(consistency, GrowthFactor, cp\_camb, likelihoods) & $\approx0.01$ \\
\midrule
\textbf{total} & \textbf{1.2} \\
\bottomrule
\end{tabular}
\end{center}
\texttt{Shear1hMisSel} is $38\%$ of the total --- the single largest
module in the pipeline, and the one this note has been analysing.
The reason it costs $4.8\times$ more than \texttt{NumCountsSel}, even
though both run the identical 2-D $(\lnM,z)$ quadrature per bin, is
exactly the obstruction of \S\ref{sec:exact}--\S\ref{sec:setup}:
\texttt{NumCountsSel} runs that quadrature once per bin (12 calls);
\texttt{Shear1hMisSel} reruns it once per \emph{(bin, radius)} pair
(12 bins $\times$ 10 radii $=120$ calls, via
\texttt{use\_cartesian\_product}), because today nothing separates
the $R$-dependence from the mass integral. Implementing Idea 2 or
Idea 3 removes exactly that repetition: the expensive quadrature
(the moments $\mu_n^{(ij)}$, or the Chebyshev projections
$c_k^{(ij)}$) is computed once per bin, the same 12 calls
\texttt{NumCountsSel} already makes, and the 120 \texttt{(bin,
radius)} evaluations become closed-form sums, not adaptive Cuhre
calls. That predicts \texttt{Shear1hMisSel} dropping from $0.458$s to
something in \texttt{NumCountsSel}'s own cost class, $\sim0.1$s ---
taking the pipeline total from $1.2$s to $\sim0.84$s, under the
$1$s target.

\section{The C++ implementation and its measured payoff}
\label{sec:cpp_impl}

The prediction above has now been implemented, measured, and shipped
\emph{in place}: \texttt{Shear1hMisSel.so} and \texttt{NumCountsSel.so}
themselves are now fixed-GL evaluators
(\pyfile{src/models/n\_operator\_sel\_gl\_t.hh}, drivers
\pyfile{src/modules/num\_counts\_sel/\{Shear1hMis,NumCounts\}.cc}), in
the same family as \texttt{b\_sel\_marg} and
\texttt{shear\_prj\_frozen\_physics}: no Cuhre calls at all. A shared
core performs one fixed Gauss--Legendre $(z,\lnM)$ pass per bin per
sample (defaults $n_{\lnM}=96$, $n_z=64$ nodes) to build the exact
z-marginalised weight $\widehat W_{ij}(\lnM)$ of
eq.~\eqref{eq:reduction}. For \texttt{NumCountsSel} the observable
\emph{is} the weight's norm ($f=1$), so the module reduces to reading
off $\int d\lnM\,\widehat W_{ij}$; for \texttt{Shear1hMisSel} the
weight additionally folds in the $z$-only
$\Sigma_{\rm crit}^{-1}(z)$ factor and every radius is served from the
cached weight. Two service methods are ini-selectable for the shear
module (\texttt{method = exact | idea2}): \emph{exact} (the default)
evaluates the remaining 1-D mass sum
$\sum_k w_k \widehat W_{ij,k}\,\Phi_i(R,\lnM_k)$ per (bin, $R$) point;
\emph{idea2} uses eq.~\eqref{eq:approach2_general} at $n\le2$ with a
3-point central-difference stencil for $\Phi''$. Module labels, grid
semantics (\texttt{bin\_index} wall; \texttt{bin\_index} $\times$
\texttt{r\_perp} cartesian product, bin-slow/$R$-fast) and output
sections (\texttt{numcountssel/vals}, \texttt{shear1hmissel/vals}) are
unchanged, so the production ini and
\pyfile{y3\_buzzard/likelihood\_cp.py} work as-is --- the retired
adaptive-integrator knobs (\texttt{algorithm}, \texttt{eps\_rel},
\texttt{max\_eval}, \dots) are simply ignored.\footnote{One
implementation note: the output section names are deliberately
hardcoded rather than ini-configurable. A CosmoSIS \texttt{[DEFAULT]}
block propagates its keys into \emph{every} module section, so an
\texttt{output\_section} knob silently redirected the modules' output
to the vestigial \texttt{cluster\_observables\_full} value in the
production ini's \texttt{[DEFAULT]} --- found the hard way during
validation.}

Beyond the mean speed-up, fixed GL makes the cost
\emph{deterministic}. Adaptive Cuhre's runtime depends strongly on the
sampled parameter point: over $\sim10^6$ MCMC realisations of the
production chain, \texttt{NumCountsSel} averaged $0.107$\,s but
ranged up to $0.98$\,s ($9\times$ its mean), and
\texttt{Shear1hMisSel} averaged $0.575$\,s with a $4.0$\,s tail
($7\times$) --- the integrator running wild trying to converge at hard
corners of parameter space. The GL evaluators execute an identical
node count at every sample, so those tails vanish by construction.

Measured on the extraction pipeline
(\pyfile{docs/figs/real\_pipeline\_extract.ini}, login node, full
\texttt{bin\_index} 0--11), against the retired Cuhre modules:
\begin{center}
\small
\begin{tabular}{l r r}
\toprule
 & Cuhre (old) [s] & fixed-GL (new) [s] \\
\midrule
\texttt{NumCountsSel}   & 0.088 & \textbf{0.021} \\
\texttt{Shear1hMisSel}  & 0.470 & \textbf{0.028} \\
\bottomrule
\end{tabular}
\end{center}
(\texttt{NumCountsSel} agrees with its Cuhre predecessor to
$3.1\times10^{-4}$ across all 12 bins, inside the old module's own
\texttt{eps\_rel}$=1.5\times10^{-3}$.) The full production pipeline
(\texttt{mock\_mcmc\_cp\_camb.ini}, \texttt{test} sampler, likelihood
included) now completes in \textbf{0.70\,s} per sample, down from
1.2\,s --- under the 1\,s target, with the two converted modules
together costing $7\%$ of the budget instead of $46\%$. \texttt{idea2}
buys only a few ms more than \texttt{exact} while adding its
truncation error, confirming \S\ref{sec:comparison}'s expectation that
once the z-marginalisation is in place the exact 1-D sum is already
essentially free; \texttt{exact} is the production default.

Accuracy, against the old module's own output:
\begin{itemize}
\item \textbf{Bins 0--3}: $\max|{\rm new}/{\rm old}-1| = 3.3\times10^{-4}$
  --- an order of magnitude \emph{inside} the old module's own Cuhre
  tolerance (\texttt{eps\_rel}$=5\times10^{-3}$).
\item \textbf{\texttt{idea2} vs \texttt{exact}}: median
  $8\times10^{-4}$, max $4.6\times10^{-3}$ --- matching
  \S\ref{sec:real_validation}'s Python-side prediction for the $n\le2$
  truncation.
\item \textbf{Quadrature convergence}: doubling to
  $(n_{\lnM},n_z)=(192,128)$ moves \texttt{exact} by
  $\le4\times10^{-4}$. The $z$ axis is the one that needs the nodes:
  $S_{ij}(\lnM,z)$ has compact support ($\Delta z\sim0.15$ inside the
  full $[0.05,0.80]$ window), which is why the default is $n_z=64$
  rather than 32.
\end{itemize}

\paragraph{One deliberate physics change: bins 4--11.} The old
module's \texttt{set\_bin} is a silent no-op for
\texttt{bin\_index}$\ge4$ (\texttt{lob\_centers\_} holds only the 4
DES-Y3 richness centres), so bins 4--11 --- the $z\in[0.35,0.50]$ and
$[0.50,0.65]$ blocks --- were reusing bin 3's
$R_\lambda(130)$ in the miscentered term instead of their own richness
bin's aperture. The new \texttt{Shear1hMisSel} resolves the richness bin as
\texttt{bin\_index} $\bmod\ 4$ (the \texttt{sel\_function} layout is
richness-fast, $z$-block-slow), which changes bins 4--11 by up to
$2.2\%$ (largest for richness bin 0, where
$R_\lambda(25)/R_\lambda(130)$ differs most). Two checks pin this down
as the entire difference: (i) bins 7 and 11 --- richness bin 3, whose
$R_\lambda(130)$ the old module was accidentally using anyway ---
agree with the old module at the same $\lesssim10^{-3}$ level as bins
0--3; and (ii) forcing \texttt{lob\_centers = 130 130 130 130} in the
new module reproduces the old module's bins 4--11 to
$1.6\times10^{-3}$. The new values are the correct ones.

\end{document}
